% =====================================================================
% Section 6 — Evaluation
% Target: Software: Practice and Experience (Wiley). Style-agnostic:
% requires booktabs, siunitx, and (optionally) threeparttable.
% Numbers are the official native-Linux run (see paper/tables/results.tex).
% Regenerate data from: benchmarks/results/{RESULTS-docker.md,STATS.md,raw-docker.csv}
% =====================================================================
\section{Evaluation}
\label{sec:eval}

We evaluate JxMVC against four widely used JVM web frameworks---Spring
Boot, Quarkus, Micronaut, and Javalin---plus a GraalVM native-image
build of Quarkus, under identical, containerised conditions. The
evaluation is organised around three research questions:

\begin{description}
  \item[RQ1 (Footprint).] How does JxMVC compare on deployable image
    size, cold-start time, and resident memory?
  \item[RQ2 (Performance).] How does JxMVC compare on request
    throughput and tail latency for equivalent endpoints?
  \item[RQ3 (Correctness under load).] Does JxMVC serve sustained
    concurrent load without errors or non-2xx responses?
\end{description}

\noindent The cost of security-by-default (RQ4 in our design notes) is
discussed qualitatively in Section~\ref{sec:threats}, as it is a
property of the request pipeline rather than a throughput dimension.

\subsection{Methodology}
\label{sec:eval-method}

\paragraph{Harness.} All frameworks are built and measured by a single
reproducible, one-command Docker harness that ships with the artifact
(\texttt{benchmarks/}). Each framework is packaged as an image on the
\emph{same} JRE base, run under \emph{identical} CPU and memory limits,
and exposes two functionally equivalent endpoints: \texttt{/plaintext}
(a fixed string) and \texttt{/json} (a small serialised object). The
harness measures deployable image size, cold-start time (from
\texttt{docker run} to the first \texttt{HTTP 200} on \texttt{/plaintext}),
resident set size (RSS, via \texttt{docker stats} under load), and
throughput/latency generated by a dependency-free JDK load client.

\paragraph{Protocol.} Following established guidance on rigorous JVM
performance measurement~\cite{georges2007,kalibera2013}, we use $C=64$
concurrent connections, a $30$\,s measurement window, and $R=5$
repetitions per (framework, endpoint) cell, each preceded by a $5$\,s
warm-up phase so that the JIT reaches steady state before measurement.
We report the \emph{median} across repetitions and, for throughput, the
observed $[\min,\max]$ range.

\paragraph{Isolation and reproducibility.} To prevent the load
generator from contending with the server under test, the container is
pinned to two physical performance cores
(\texttt{--cpuset-cpus=0-3}, \texttt{--cpus=4}, \texttt{--memory=2g})
and the load client is pinned to two \emph{disjoint} performance cores
(\texttt{taskset -c 4-7}); operating-system and daemon work is left to
the remaining cores. The CPU governor is fixed to \texttt{performance}
and Turbo Boost is disabled so that clock frequency does not vary with
temperature across the run. Table~\ref{tab:env} reports the full
environment.

\begin{table}[t]
  \centering
  \caption{Evaluation environment.}
  \label{tab:env}
  \small
  \begin{tabular}{@{}ll@{}}
    \toprule
    Component & Value \\
    \midrule
    CPU        & Intel Core i5-12500H (4 P-cores @2.5\,GHz, 8 E-cores @1.8\,GHz) \\
    Memory     & 30\,GiB \\
    Kernel     & Linux 7.0.14 (bare-metal Arch Linux) \\
    Host JDK   & OpenJDK 25 (load client) \\
    Container  & Docker 29.6.1, shared JRE base image \\
    CPU state  & governor \texttt{performance}, Turbo Boost off \\
    Container limits & \texttt{--cpus=4 --memory=2g --cpuset-cpus=0-3} \\
    Load client & pinned to cores 4--7 (\texttt{taskset}) \\
    Load        & $C=64$, $30$\,s $\times\,5$ reps, $5$\,s warm-up \\
    \bottomrule
  \end{tabular}
\end{table}

\subsection{Results}
\label{sec:eval-results}

Table~\ref{tab:results} reports the headline metrics and
Table~\ref{tab:latency} the per-request latency. Across the entire
campaign the frameworks served \num{92759575} requests
($\approx$\,92.8\,million) with \textbf{zero} errors and zero non-2xx
responses, directly answering RQ3.

% Main results table — median [min–max] over 5 reps.
% Source: benchmarks/results/STATS.md (native-Linux run, i5-12500H).
% Requires: booktabs, siunitx. Best value per column in \textbf.
\begin{table*}[t]
  \centering
  \caption{Headline benchmark results. Image size, cold start, and RSS
    are per framework; throughput is the median over five 30\,s
    repetitions with the observed $[\min,\max]$ range. Load: 64
    concurrent connections. Best JVM-mode value per column in bold;
    the GraalVM native build is listed separately for reference.}
  \label{tab:results}
  \sisetup{group-separator={,}}
  \begin{tabular}{@{}l S[table-format=3.1] S[table-format=4.0] S[table-format=3.1] r r@{}}
    \toprule
    & {Image} & {Startup} & {RSS} & {\texttt{/plaintext}} & {\texttt{/json}} \\
    Framework & {(MB)} & {(ms)} & {(MB)} & {req/s (med [min--max])} & {req/s (med [min--max])} \\
    \midrule
    JxMVC      & \bfseries 271.7 & 822  & 448.5          & 49\,062 [47\,198--49\,645] & 48\,720 [48\,389--49\,005] \\
    Spring Boot& 299.9           & 1945 & 375.6          & 49\,462 [47\,865--49\,979] & 49\,896 [48\,789--50\,295] \\
    Quarkus    & 295.9           & 606  & 431.2          & \bfseries 55\,087 [54\,767--55\,552] & \bfseries 54\,209 [53\,831--55\,391] \\
    Micronaut  & 292.5           & 1154 & \bfseries 331.4& 50\,974 [50\,502--51\,922] & 50\,652 [50\,295--51\,444] \\
    Javalin    & 286.4           & \bfseries 369 & 424.6 & 54\,169 [53\,567--57\,131] & 53\,614 [52\,926--54\,148] \\
    \midrule
    \multicolumn{6}{@{}l}{\itshape GraalVM native (Quarkus), for reference} \\
    Quarkus-native & 72.8 & 12 & 25.1 & 52\,699 [52\,488--53\,123] & 50\,073 [46\,614--52\,214] \\
    \bottomrule
  \end{tabular}
\end{table*}

% Latency table — median across reps of per-rep mean and p99 (ms).
% Source: benchmarks/results/raw-docker.csv (cols meanMs, p99).
% Requires: booktabs, siunitx.
\begin{table}[t]
  \centering
  \caption{Per-request latency (milliseconds): mean and 99th
    percentile, median across repetitions, at 64 concurrent
    connections. All JVM frameworks fall within a narrow band.}
  \label{tab:latency}
  \small
  \begin{tabular}{@{}l S[table-format=1.2] S[table-format=1.2] S[table-format=1.2] S[table-format=1.2]@{}}
    \toprule
    & \multicolumn{2}{c}{\texttt{/plaintext}} & \multicolumn{2}{c}{\texttt{/json}} \\
    \cmidrule(lr){2-3}\cmidrule(lr){4-5}
    Framework & {mean} & {$p_{99}$} & {mean} & {$p_{99}$} \\
    \midrule
    JxMVC          & 1.30 & 3.64 & 1.31 & 3.59 \\
    Spring Boot    & 1.29 & 3.31 & 1.28 & 3.29 \\
    Quarkus        & 1.16 & 3.19 & 1.18 & 3.12 \\
    Micronaut      & 1.25 & 3.16 & 1.26 & 3.56 \\
    Javalin        & 1.18 & 3.31 & 1.19 & 3.40 \\
    \midrule
    Quarkus-native & 1.21 & 3.64 & 1.28 & 3.85 \\
    \bottomrule
  \end{tabular}
\end{table}


\subsection{Analysis}
\label{sec:eval-analysis}

\paragraph{RQ1 --- Footprint.} JxMVC produces the \emph{smallest
JVM-mode deployable image} of the group (\SI{271.7}{\mega\byte}, about
5\% smaller than the next-smallest, Javalin at
\SI{286.4}{\mega\byte}, and \SI{28}{\mega\byte} below Spring Boot).
Its cold start (\SI{822}{\milli\second}) is $2.4\times$ faster than
Spring Boot (\SI{1945}{\milli\second}) and faster than Micronaut
(\SI{1154}{\milli\second}), though slower than Quarkus
(\SI{606}{\milli\second}) and Javalin (\SI{369}{\milli\second}).
The clear cost is resident memory: JxMVC's RSS
(\SI{448.5}{\mega\byte}) is the highest of the JVM stacks, a direct
consequence of running on an embedded Tomcat servlet container. We do
not hide this; it is the principal footprint trade-off of the design
and is revisited in Section~\ref{sec:threats}. For context, the
GraalVM native build reduces all three dimensions by roughly an order
of magnitude (\SI{72.8}{\mega\byte} image, \SI{12}{\milli\second}
start, \SI{25.1}{\mega\byte} RSS), quantifying what JxMVC would gain
from future ahead-of-time compilation.

\paragraph{RQ2 --- Performance.} On throughput, JxMVC is
\emph{competitive but not the leader}. It sustains
\SI{49062}{req\per\second} on \texttt{/plaintext} and
\SI{48720}{req\per\second} on \texttt{/json}, statistically on par
with Spring Boot (within $\sim$1\%) and within roughly 4\% of
Micronaut. Quarkus (\SI{55087}{req\per\second}) and Javalin
(\SI{54169}{req\per\second}) lead by $\sim$11\% and $\sim$9\%
respectively. Latency tells a tighter story: mean service time is
\SI{1.30}{\milli\second} and the 99th percentile
\SI{3.6}{\milli\second}, essentially indistinguishable from the rest
of the field (all frameworks fall within \SIrange{3.1}{3.9}{\milli\second}
at $p_{99}$). In other words, a zero-dependency stack built directly
on the JDK and Jakarta Servlet API pays no meaningful latency penalty
and only a modest throughput gap versus mature, optimised runtimes.

\paragraph{RQ3 --- Correctness under load.} No framework, including
JxMVC, produced a single error or non-2xx response across
\num{92759575} requests. JxMVC's secure-by-default pipeline---which
adds security headers, CSRF handling, rate-limit bookkeeping, and body
caps to \emph{every} request---therefore imposes its guarantees
without sacrificing reliability under sustained concurrency.

\paragraph{Summary.} JxMVC delivers the smallest JVM image and a
faster cold start than half the field, at competitive throughput and
indistinguishable tail latency, in \SI{253}{\kilo\byte} of
zero-dependency code---its one clear trade-off being higher resident
memory from the servlet container. This positions JxMVC as a credible
choice where deployable size, auditability, and security-by-default
matter more than peak requests-per-second or minimal RAM.
